\section{Related Works}
\label{sec:related_work}

The development of \ours aligns with the rapidly evolving landscape of frameworks that leverage large language models (LLMs) for the creation of language agents and multi-agent systems. Here, we briefly introduce works closely related to \ours from two sub-domains pertinent : Language Agent Frameworks, focusing on individual agent capabilities, and Multi-Agent Frameworks, emphasizing collaboration among multiple agents. For broader related works, readers can refer to \citep{agent-survey-1,agent-survey-2}.

\paragraph{Language Agent Frameworks}
Language agent frameworks are pivotal for developing applications that can interpret and interact using human language. 

The Transformers library \citep{transformers-agent} has introduced a natural language API to interface with transformer models in its recent updates (\textit{Transformers-Agents}). This API utilizes a set of customizable tools, allowing the model to interpret instructions and generate code snippets accordingly. It offers support for various open-source and proprietary model endpoints, catering to diverse developer needs.
\textit{LangChain} \citep{langchain} provides a framework for building applications that are context-aware and capable of reasoning. It includes libraries and templates that facilitate the integration of multiple components into a unified cognitive architecture. LangServe and LangSmith extend the framework's capabilities by enabling deployment as a REST API and offering developer tools for debugging and monitoring chains built on any LLM framework.
\textit{AutoGPT} \citep{autogpt} illustrates a different approach, allowing an LLM to iteratively execute actions and make decisions. As a generalist agent, AutoGPT is not task-specific; it is designed to perform a variety of computer-based tasks, reflecting the adaptive nature of LLMs.
\textit{ModelScope-Agent} \citep{modelscope-agent} is a customizable agent framework that harnesses open-source LLMs to perform tasks and connect with external APIs. It facilitates seamless integration with model APIs and common APIs while providing a comprehensive infrastructure for data collection, tool retrieval, and customized model training, all aiming to realize practical real-world applications.


\paragraph{Multi-Agent Frameworks}
Building on the capabilities of individual agents, multi-agent frameworks explore collaboration and interaction among multiple agents to address complex tasks.

\textit{AutoGen} \citep{auto-gen} provides a generic infrastructure that allows developers to program interaction patterns using both natural language and code. This framework enables the development of diverse applications by facilitating conversation among agents that are customizable and can utilize various combinations of LLMs, human inputs, and tools.
\textit{MetaGPT} \citep{meta-gpt} incorporates meta-programming to enhance multi-agent collaborations. By encoding Standardized Operating Procedures (SOP) into prompts, this framework ensures streamlined workflows and reduced errors, exemplifying effective task decomposition among agents.
\textit{AGENTS} \citep{Agents} is an open-source library that supports autonomous language agents with features like planning, memory, and multi-agent communication. It is designed to be user-friendly, helping non-specialists to deploy state-of-the-art language agents, and research-friendly, with a modularized design for extensibility.
\textit{OpenAgents} \citep{OpenAgents} provides an open platform for using language agents with practical functionalities accessible through a web interface. This framework emphasizes facilitating real-world agent interactions and includes specialized agents for different tasks, such as data analysis and web browsing.
\textit{ChatDev} \citep{ChatDev} exploits LLMs for software development, creating a virtual chat-powered company that follows a waterfall model. It engages ``software agents'' at different stages of the development process, facilitating collaboration and context-aware communication.
\textit{CAMEL} \citep{camel} proposes a novel framework for autonomous cooperation among communicative agents using role-playing techniques, which allows for the generation of conversational data for studying agent behaviors and capabilities.
Lastly, \textit{AgentSims} \citep{AgentSims} introduces a sandbox environment to evaluate LLMs in task-based scenarios, offering an infrastructure for researchers to test specific LLM capacities within a simulated environment.

These frameworks represent significant strides in the use of LLMs for both individual and collaborative agent tasks. \ours is situated within this context, contributing by addressing the need for a user-friendly, fault-tolerant and versatile framework designed to manage complex interactions and processes inherent in multi-agent LLM systems. By focusing on ease of use and reliability, \ours aims to facilitate the creation of robust and versatile applications across diverse domains.
